% Every pointer was grepped in book/chapters and book/frontmatter. LNO and MGNO occur
% nowhere in the book and are recorded as absent.
% Bibliography pass: the two tomographic and the one tensor-factorized variant rows were
% deleted together with their arXiv-only entries; multiplier-array factorization survives
% above as a technique remark tied to the book's own \eqref{eq:ch8-tucker}, carrying no
% citation and no measured compression ratio.
\chapter{Index of Architectural Variants}\label{app:e}
\index{architectural variants!index of|textbf}
\index{LNO (linear neural operator)!not treated in this book}
\index{MGNO (multi-grid neural operator)!not treated in this book}
\index{FourCastNet|see{adaptive Fourier neural operator (AFNO)}}

Table~\ref{tab:synthesis_matrix} of
\chapref{ch:part3-synthesis-matrix-and-real-world-deployment} orders the deployed
architectures by the domain that consumes them. This appendix inverts that ordering for
the reader who arrives with a name and wants the place the book builds it, the citation
attached there, and the one constraint separating the variant from the layer it modifies.
Nothing here is new: each distinguishing feature restates a restriction already written
out in \S\ref{sec:ch8-axes}, where every architecture is the generic integral-kernel layer
\eqref{eq:no_layer} with the kernel's algebraic form, its support, or the representation of
its input and output constrained. Which variants carry a theorem is kept in
Table~\ref{tab:ch8-ledger} and is not duplicated below. Section~\ref{sec:ape-notation}
then turns the front matter's Table~\ref{tab:notation-master} into a locator for the
argument.

\section{The Variants by Name}
\label{sec:ape-variants}
% -----------------------------------------------------------------------------

Read Table~\ref{tab:variant-index} as a set of deltas from one template. Its
\emph{introduced in} column names where the variant is built or defined, not where its
name first appears: \chapref{ch:part1-introduction-the-paradigm-shift-to-infinit}
previews the graph, geometry-informed, physics-informed, multi-scale and adaptive-Fourier
descendants of the spectral layer in a single passage, and
\S\ref{sec:ch8-requirements} states their hypotheses and costs, but the construction each
citation points to lies in Part~II or Part~III. The conclusion the table bears out is the
taxonomy chapter's: these variants differ by \emph{where} they spend their restriction, and
none is a rival to the Fourier neural operator (FNO) of \eqref{eq:fno_layer} as a general
method. The FNO, the generic neural operator and DeepONet are primary objects of
\chapref{ch:part1-the-fourier-neural-operator-spectral-theor},
\chapref{ch:part1-operator-approximation-theory-and-universa} and
\chapref{ch:part2-deeponet-engineering-branch-trunk-networks}, and are not indexed here as
variants. Neither is the low-rank multilinear reconstruction of the multiplier array: that is
a factorization technique acting on the weights of an existing layer, its arithmetic is
\eqref{eq:ch8-tucker} in \S\ref{sec:ch8-cost}, what it reduces is parameter count, and the
premise that the multipliers of a given solution operator are genuinely low-rank is empirical
rather than proved.

\begin{table}[htbp]
\centering
\footnotesize
\caption{\textbf{Index of architectural variants}, read from the architecture side. Names,
citations and constraints are as printed in the chapter named, compressing the analysis of
\S\ref{sec:ch8-axes}.}
\label{tab:variant-index}
\begin{tabularx}{\textwidth}{@{}l >{\raggedright\arraybackslash}X l
  >{\raggedright\arraybackslash}X >{\raggedright\arraybackslash}X@{}}
\toprule
\textbf{Variant} & \textbf{Name as printed} & \textbf{Cited} & \textbf{Introduced in} & \textbf{Distinguishing constraint} \\
\midrule
GNO & graph neural operator & \cite{li2023geometry,brandstetter2022message} &
\S\ref{sec:ch24-continuous-kernel} of \chapref{ch:part2-graph-neural-operators-gno-and-irregular-g} &
Support restriction: the layer \eqref{eq:no_layer} integrates over the ball
$\mathcal{B}_r(x_i)$ of a mesh, with Voronoi weights. \\
GINO & geometry-informed neural operator & \cite{li2023geometry} &
\S\ref{sec:ch24-gino}, same chapter; applied in \S\ref{sec:ch32-gino} &
Graph encoder and decoder bracket a latent-grid FNO, as composed in \eqref{eq:ch8-gino}. \\
AFNO & adaptive Fourier neural operator & \cite{guibas2021afno,pathak2022fourcastnet} &
\S\ref{sec:ch33-fourcastnet} of \chapref{ch:part3-global-weather-forecasting-and-climate-dyn} &
Patch token mixer; soft thresholding \eqref{eq:ch33-shrink} makes retained frequencies
input-dependent, not a fixed box \eqref{eq:ch05-modeset}. \\
U-FNO & U-Net-enhanced Fourier neural operator & \cite{wen2022ufno} &
\chapref{ch:part3-subsurface-multiphase-flow-and-carbon-stor}, its U-FNO
section$^{\dagger}$ &
Adds a parallel local branch $\mathcal{U}_\ell$ to the spectral layer \eqref{eq:ch34-layer},
for fronts whose spectral tail is \eqref{eq:ch09-jump-tail}. \\
PINO & physics-informed neural operator & \cite{li2021physics} &
\chapref{ch:part1-physics-informed-operators-and-sobolev-reg}$^{\dagger}$; applied in
\S\ref{sec:ch35-map}, \S\ref{sec:ch37-operator} &
Not a kernel restriction but the training functional \eqref{eq:pino_loss}; residual
differentiation exact on the torus via \eqref{eq:spectral_diff}, approximate elsewhere. \\
\bottomrule
\end{tabularx}
\end{table}

Three readings need guarding. First, the dagger marks the two loci that resolve at chapter
granularity because the sections concerned carry no \texttt{sec:} label --- the U-FNO
section of the subsurface chapter, and the physics-informed chapter entire; a label there
would sharpen the pointer without changing what it names. Second, the book attaches two
names at different levels in \S\ref{sec:ch33-fourcastnet}: FourCastNet is the weather
system, AFNO \cite{pathak2022fourcastnet,guibas2021afno} the mixer it is built from, and the
synthesis matrix's weather row names the layer --- one architecture, one row here. Last, two
variants that would round out the index are missing from it: neither the linear neural
operator (LNO) nor the multi-grid neural operator (MGNO) is introduced, cited or mentioned
anywhere in this book, and no bibliographic entry corresponds to either. The nearest
material is negative --- \chapref{ch:part3-the-multi-query-bottleneck-in-computationa}
names multigrid only as a feature of the classical solver being accelerated --- so both are
indexed as absent rather than left to be discovered by omission.

\section{Symbols to Their Chapters}
\label{sec:ape-notation}
% -----------------------------------------------------------------------------

The master notation table defines each symbol where it is defined; it cannot say where the
symbol begins to carry weight, which is what a reader returning to an argument needs.
Table~\ref{tab:symbol-map} gives one pointer per principal symbol, anchored on an equation
or section label so it can be followed mechanically. The conclusion is that the book's
notation is local to its parts: the objects of \eqref{eq:fredholm_kernel} are never
restated when Part~II writes them as tensor axes, and Part~III introduces its physical
symbols with their governing equations rather than with the network.

{\small
\begin{xltabular}{\textwidth}{@{}l >{\raggedright\arraybackslash}X@{}}
\caption{\textbf{Symbol-to-chapter map} for the principal symbols of
Table~\ref{tab:notation-master}.}
\label{tab:symbol-map}\\
\toprule
\textbf{Symbol} & \textbf{First used seriously} \\
\midrule
\endfirsthead
\multicolumn{2}{r}{\textit{Table~\ref{tab:symbol-map}, continued}}\\
\toprule
\textbf{Symbol} & \textbf{First used seriously} \\
\midrule
\endhead
\midrule
\multicolumn{2}{r}{\textit{Table~\ref{tab:symbol-map}, continued on the next page}}\\
\endfoot
\bottomrule
\endlastfoot
\multicolumn{2}{l}{\textit{Part I: the objects}}\\
\addlinespace
$\mathcal{A}$, $\mathcal{U}$, $\mathcal{G}_\theta$, $\mu$ & \eqref{eq:ch01-solution-operator}
and the risk \eqref{eq:ch01-risk}, where $\mu$ is the input measure. \\
$L^{p}$, $H^{s}$, $D^{\alpha}u$, $\mathcal{L}(X,Y)$, $\langle\cdot,\cdot\rangle_{X^{*},X}$ &
\chapref{ch:part1-function-spaces-duality-and-sobolev-regula}, via
Theorem~\ref{thm:ch02-riesz-rep} and the compact embedding Theorem~\ref{thm:ch02-rellich}. \\
$\mathbb{T}^{d}$, $\mathcal{F}$, $\mathcal{F}^{-1}$, $R_{\phi}(k)$ & The torus as evolution
domain in \chapref{ch:part1-partial-differential-equations-as-solution}; as the FNO's
hypothesis, and the transform as multiplier calculus, in \S\ref{sec:ch05-multipliers}. \\
$Z_{k_{\max}}$, $k_{\max}$ & \eqref{eq:ch05-modeset} in \S\ref{sec:ch05-truncation}. \\
$\kappa_{\ell}$, $W_{\ell}$, $d_v$, $L$ & \eqref{eq:no_layer} and \eqref{eq:fredholm_kernel}
of Definition~\ref{def:generic_no}. \\
$\mathcal{M}_N$, $P_N$, $h_N$, $w_j$, $\Delta x$ & \S\ref{sec:ch06-definition} and
\S\ref{sec:ch06-quadrature}; the weights become mesh-derived in
\S\ref{sec:ch24-continuous-kernel}. \\
$\mathcal{N}(\cdot\,;a)$, $\mathcal{B}$, $\lambda_{\text{bc}}$ & \eqref{eq:pino_loss}; their
dimensionality is restated for boundaries in \S\ref{sec:ch37-operator}. \\
\addlinespace
\multicolumn{2}{l}{\textit{Part II: the tensors}}\\
\addlinespace
$B$, $C$, $C_{\mathrm{in}}$, $C_{\mathrm{out}}$, $H$, $W$, $X_{\mathrm{nodes}}$, $N_{\mathrm{nodes}}$ &
\eqref{eq:ch21-grid-layout} and \eqref{eq:ch21-node-layout} in \S\ref{sec:ch21-layouts};
coordinate augmentation in \eqref{eq:ch21-augment}. \\
$m$, $P$, $p$, $h$ & DeepONet's sensors, queries, latent and hidden widths in
\S\ref{sec:ch22-blueprint}. \\
$m_x$, $m_y$, $Z_m$, $R^{(1)}$, $R^{(2)}$ & \eqref{eq:ch23-corners} and
\eqref{eq:ch23-Rtensors} in \S\ref{sec:ch23-spectral-arithmetic}. \\
$M$, $\mu(x)$, $\sigma(x)$, $r$, $\mathcal{N}_r(i)$ & \eqref{eq:ch21-normalize} in
\S\ref{sec:ch21-normalization}; the radius $r$ and its ball in
\S\ref{sec:ch24-continuous-kernel}. \\
\addlinespace
\multicolumn{2}{l}{\textit{Part III: the physics}}\\
\addlinespace
$Q$, $Q^{\star}$, $\mathcal{T}$, $\Delta_{\mathcal{T}}$ & \eqref{eq:ch31-ledger} in
\S\ref{sec:ch31-hypothesis}; crossover \eqref{eq:ch38-crossover} in \S\ref{sec:ch38-columns};
conservation defect \eqref{eq:ch38-defect} in \S\ref{sec:ch38-roadblocks}. \\
$Re$, $U_\infty$, $C_d$, $C_l$, $y^{+}$, $\tau_w$, $\mu_t$ & \eqref{eq:ch32-reynolds} and
\eqref{eq:ch32-force-coeff} in \S\ref{sec:ch32-formulation}. \\
$\mathbf{u}_h$, $w$, $\theta$, $f$, $\Omega_{\mathrm{rot}}$, $\phi$ &
\S\ref{sec:ch33-primitive}, where $\phi$ is latitude and $w$ vertical wind. \\
$\mathbf{K}(x)$, $\phi(x)$, $S_w$, $S_g$, $k_{r\alpha}$, $\mu_\alpha$ &
\eqref{eq:ch34-cons-w} and \eqref{eq:ch34-flux} of
\chapref{ch:part3-subsurface-multiphase-flow-and-carbon-stor}. \\
$V_m$, $C_m$, $\chi$, $\mathbf{D}(x)$, $\mathbf{w}$, $h$, $\Delta t$ & \eqref{eq:monodomain},
\eqref{eq:ch35-recovery} and \eqref{eq:ch35-cost} of
\chapref{ch:part3-cardiac-electrophysiology-and-arrhythmia-d}, where $h$ is micrometers. \\
$\mathcal{R}$, $\mathbf{M}$, $\mathcal{A}$, $p(\theta,s)$, $m_k$ & \eqref{eq:ch36-radon} in
\S\ref{sec:ch36-radon}; the sampled operator recurs in \eqref{eq:ch36-consistency}. \\
$\mathrm{ACC}$, $\mathrm{RQE}$, $\mathrm{PSNR}$, $\mathrm{SSIM}$, $\mathrm{RMSE}$ &
\eqref{eq:ch33-acc} in \S\ref{sec:ch33-verification}, \eqref{eq:ch33-rqe} in
\S\ref{sec:ch33-limits}, \eqref{eq:ch36-metrics} in \S\ref{sec:ch36-metrics}. \\
\end{xltabular}
}

The double readings the master table warns about are visible here as pairs of loci: $h$ a
hidden width in \S\ref{sec:ch22-blueprint} and a mesh width in \eqref{eq:ch35-cost}, $M$ an
ensemble size in \eqref{eq:ch21-normalize} and a view count in \S\ref{sec:ch36-radon}, $\mu$
a measure, a pointwise mean and a phase viscosity in three rows above, $\phi$ a latitude and
a porosity in two. Reading the paired rows is what disambiguates them; the Notation
table lists the reuses a reader must actively track, and the glossary entries state
each reading where it recurs.

What this index cannot do is rank the variants, and so it carries no cost column: the
arithmetic is in \eqref{eq:ch8-costs}, the hypotheses in \S\ref{sec:ch8-requirements}, and
the measured quotients in Table~\ref{tab:synthesis_matrix}, whose baselines are choices. For definitions of the terms both tables use, see the
glossary in Appendix~\ref{app:glossary}; for page-level lookup of these architectures and
symbols, the subject index that follows carries the entries authored chapter by chapter.
